\documentclass[11pt,a4paper]{article}
% =====================================================================
%  Gemeinsame Vorlage der Arbeitsblätter – 2. Funktionsweise eines Rechners (Q13)
%  Kopf, Fuß und Boxen wie informatik/13/formale-sprachen/arbeitsblaetter/ab-vorlage.tex,
%  Antwortplatz-Makros wie informatik/10/java/arbeitsblaetter/ab-vorlage.tex
%  (Linien 8 mm, Karo 5 mm aus ganzen Kästchen, Lücken seit 2026-10 x 1,5).
%  Einbinden in jedem Blatt direkt nach \documentclass: % =====================================================================
%  Gemeinsame Vorlage der Arbeitsblätter – 2. Funktionsweise eines Rechners (Q13)
%  Kopf, Fuß und Boxen wie informatik/13/formale-sprachen/arbeitsblaetter/ab-vorlage.tex,
%  Antwortplatz-Makros wie informatik/10/java/arbeitsblaetter/ab-vorlage.tex
%  (Linien 8 mm, Karo 5 mm aus ganzen Kästchen, Lücken seit 2026-10 x 1,5).
%  Einbinden in jedem Blatt direkt nach \documentclass: % =====================================================================
%  Gemeinsame Vorlage der Arbeitsblätter – 2. Funktionsweise eines Rechners (Q13)
%  Kopf, Fuß und Boxen wie informatik/13/formale-sprachen/arbeitsblaetter/ab-vorlage.tex,
%  Antwortplatz-Makros wie informatik/10/java/arbeitsblaetter/ab-vorlage.tex
%  (Linien 8 mm, Karo 5 mm aus ganzen Kästchen, Lücken seit 2026-10 x 1,5).
%  Einbinden in jedem Blatt direkt nach \documentclass: \input{ab-vorlage}
%  Übersetzen mit XeLaTeX (oder Tectonic):
%
%  Blatt:   xelatex ab-2-2-registermaschine.tex
%  Lösung:  xelatex ab-2-2-registermaschine-lsg.tex
%           (Datei enthält nur: \def\mitloesung{}\input{ab-2-2-registermaschine})
%  Lösungseinträge belegen nie mehr Platz als die Lücke im Blatt.
%
%  Antwortplatz:
%   \luecke[l] / \lsg[l]{Text}      Lücke (Standard 6,75 cm, Fachbegriff 5,25 cm, Zahl 2,25 cm)
%   \restzeile / \lsgrest{Text}     Lücke bis Zeilenende      \lueckentext  Absatz mit Lücken
%   \linien{n} / \lsglinien{n}{T}   Schreiblinien 8 mm        \linienfeld{n} / \lsgfeld{n}{Code}
%   \schreibzeile                   Tabellenzelle 8 mm        \lsgtext{…}  nur in der Lösung
%   \kariert{n} / \lsgkariert{n}{TikZ}, \karobox{b}{h} / \lsgkarobox{b}{h}{TikZ}
%  Assembler:
%   asm           Listing, Zeilennummern = Adressen (ab 0); keine Leerzeilen, Marken in der Befehlszeile
%   asmluecke     dasselbe mit 1,5-fachem Zeilenabstand für Lücken (*\lsg[3.5cm]{JGE ende}*)
%   \nr{1}        Schrittnummer ① wie auf der Website und im Simulator
%  Struktogramme: \sgwenn, \sgschleife (siehe unten)
% =====================================================================
\usepackage[a4paper,left=16mm,right=16mm,top=23mm,bottom=13mm,headheight=12mm,headsep=3mm,footskip=6mm]{geometry}
\usepackage{iftex}
\ifPDFTeX
  \usepackage[T1]{fontenc}
  \usepackage[utf8]{inputenc}
  \usepackage{helvet}
  \usepackage[scaled=0.86]{beramono}
  \renewcommand{\familydefault}{\sfdefault}
\else
  \usepackage{fontspec}
  \setmainfont{texgyreheros}[Extension=.otf,UprightFont=*-regular,BoldFont=*-bold,ItalicFont=*-italic,BoldItalicFont=*-bolditalic]
  \setmonofont{DejaVuSansMono}[Extension=.ttf,UprightFont=*,BoldFont=*-Bold,ItalicFont=*-Oblique,BoldItalicFont=*-BoldOblique,Scale=0.86]
\fi
\usepackage[ngerman,shorthands=off]{babel}
\usepackage{amssymb}
\usepackage{xcolor,graphicx,tabularx,array,enumitem,fancyhdr,tikz,multicol,calc,etoolbox,pifont,xspace,listings}
\usepackage[most]{tcolorbox}
\usetikzlibrary{arrows.meta,positioning,calc}
\usepackage[hidelinks]{hyperref}
\setlength{\parindent}{0pt}
\setlength{\parskip}{3pt}
\setlist{nosep,leftmargin=*}

% ---------- Lösungsschalter ----------
\newif\ifloesung
\ifdefined\mitloesung\loesungtrue\fi
\newcommand{\lsgfarbe}{\color{red!75!black}}

% ---------- Kopf- und Fußzeile ----------
\newcommand{\lehrkraft}{Hau}
\newcommand{\themenbereich}{2. Funktionsweise eines Rechners}
\newcommand{\abnummer}{}
\pagestyle{fancy}\fancyhf{}
\renewcommand{\headrulewidth}{0.4pt}
\fancyhead[L]{\small Klasse 13\\Datum:}
\fancyhead[C]{\small Themenbereich:\\\themenbereich}
\fancyhead[R]{\small \lehrkraft\ifloesung\\{\lsgfarbe\bfseries Lösung}\fi}
\fancyfoot[C]{\scriptsize edu-mrh.de\,·\,Informatik 13\,·\,Blatt \abnummer\ifloesung\ (Lösung)\fi\,·\,Seite \thepage}
\newcommand{\abtitel}[2]{\renewcommand{\abnummer}{#1}\begin{center}\Large\bfseries #1\quad\underline{#2}\end{center}\vspace{-2mm}}

% ---------- Boxen ----------
\newenvironment{aufgabe}[2][]{\begin{tcolorbox}[enhanced,breakable,colback=white,colframe=black,boxrule=0.7pt,arc=3mm,left=2mm,right=2mm,top=1mm,bottom=1.2mm,before skip=2mm,after skip=2mm]\textbf{Aufgabe #2}\ifstrempty{#1}{}{ \textit{#1}}\par\vspace{0.3mm}}{\end{tcolorbox}}
\newtcolorbox{merke}[1][Merke]{enhanced,breakable,colback=green!5,colframe=green!40!black,boxrule=0.8pt,arc=2mm,left=2mm,right=2mm,top=1mm,bottom=1.2mm,before skip=2mm,after skip=2mm,title={#1},fonttitle=\bfseries,coltitle=white,toptitle=0.3mm,bottomtitle=0.3mm}
\newtcolorbox{hinweis}{enhanced,breakable,colback=yellow!12,colframe=orange!70!black,boxrule=0.6pt,arc=2mm,left=2mm,right=2mm,top=1mm,bottom=1mm,before skip=2mm,after skip=2mm}
\newcommand{\web}[1]{{\small\color{gray}\textit{Website: #1}}}
\newcommand{\code}[1]{\texttt{#1}}

% ---------- Antwortplatz (wie informatik/10/java, Lücken x 1,5) ----------
\newlength{\zeilenabstand}\setlength{\zeilenabstand}{8mm}
\newlength{\karo}\setlength{\karo}{5mm}
\tikzset{karo/.style={draw=gray!50,line width=0.3pt},schreiblinie/.style={draw=black!65,line width=0.35pt}}
\newlength{\lsgbreite}
\newcommand{\lsg}[2][6.75cm]{\leavevmode\ifloesung\settowidth{\lsgbreite}{\,#2\,}%
  \ifdim\lsgbreite>#1\relax\rlap{\raisebox{0.8pt}{\resizebox{#1}{!}{\lsgfarbe\,#2\,}}}%
  \else\rlap{\raisebox{0.8pt}{\lsgfarbe\,#2}}\fi\fi\rule[-0.6ex]{#1}{0.4pt}\xspace}
\newcommand{\luecke}[1][6.75cm]{\lsg[#1]{}}
\newcommand{\lsgrest}[1]{\leavevmode\unskip\hspace{1.5mm}\ifloesung\rlap{\raisebox{0.8pt}{\lsgfarbe\,#1}}\fi\leaders\hbox{\rule[-0.6ex]{0.5pt}{0.4pt}}\hfill\null\par}
\newcommand{\restzeile}{\lsgrest{}}
\newcommand{\lsgtext}[1]{\ifloesung\leavevmode{\lsgfarbe#1}\fi}
\newcommand{\lsgoder}[2]{\ifloesung\leavevmode#1\else#2\fi}
\newcommand{\lueckentext}{\linespread{1.5}\selectfont\raggedright}
\newcommand{\schreibzeile}{\rule[-2.8mm]{0pt}{8mm}}
\newcommand{\nr}[1]{\ding{\numexpr171+#1\relax}}
\newcommand{\lsglinien}[2]{\par\vspace{0.5mm}\noindent\begin{tikzpicture}
  \path (0,0) (\linewidth-0.4pt,-#1\zeilenabstand-1.5mm);
  \foreach \i in {1,...,#1}{\draw[schreiblinie] (0.2pt,-\i\zeilenabstand) -- (\linewidth-0.2pt,-\i\zeilenabstand);}
  \ifloesung\node[overlay,anchor=north west,inner sep=0pt] at (1mm,-\zeilenabstand+1pt+5mm) {\parbox[t]{\linewidth-2mm}{\lsgfarbe\raggedright\setlength{\baselineskip}{\zeilenabstand}\rule{0pt}{5mm}#2}};\fi
  \end{tikzpicture}\par}
\newcommand{\linien}[1]{\lsglinien{#1}{}}
\newcommand{\lsgfeld}[2]{\noindent\begin{tikzpicture}
  \draw[line width=0.4pt] (0.2pt,0) rectangle (\linewidth-0.2pt,-#1\zeilenabstand-3mm);
  \foreach \i in {1,...,#1}{\draw[schreiblinie] (2mm,-\i\zeilenabstand) -- (\linewidth-2mm,-\i\zeilenabstand);}
  \ifloesung\node[overlay,anchor=north west,inner sep=0pt] at (2.5mm,-\zeilenabstand+1pt+5mm) {\parbox[t]{\linewidth-5mm}{\lsgfarbe\ttfamily\small\setlength{\baselineskip}{\zeilenabstand}\rule{0pt}{5mm}#2}};\fi
  \end{tikzpicture}}
\newcommand{\linienfeld}[1]{\lsgfeld{#1}{}}
\newcommand{\karoraster}[3]{\begin{tikzpicture}\draw[karo] (0,0) grid[step=\karo] (#1\karo,#2\karo);
  \ifloesung\begin{scope}[overlay,shift={(0,#2\karo)},red!75!black,text=red!75!black,line width=0.8pt]#3\end{scope}\fi\end{tikzpicture}}
\newcommand{\lsgkariert}[2]{\par\vspace{1.5mm}\noindent\pgfmathtruncatemacro{\karoX}{floor(\linewidth/\karo)}%
  \makebox[\linewidth]{\karoraster{\karoX}{#1}{#2}}\par\vspace{1mm}}
\newcommand{\kariert}[1]{\lsgkariert{#1}{}}
\newcommand{\lsgkarobox}[3]{\pgfmathtruncatemacro{\karoX}{round((#1)/\karo)}\pgfmathtruncatemacro{\karoY}{round((#2)/\karo)}\karoraster{\karoX}{\karoY}{#3}}
\newcommand{\karobox}[2]{\lsgkarobox{#1}{#2}{}}
% Ankreuzfeld: \kreuz{1} = angekreuzt (nur in der Lösung)
\newcommand{\kreuz}[1]{\leavevmode\ifloesung\ifnum#1=1 {\lsgfarbe$\boxtimes$}\else$\square$\fi\else$\square$\fi}

% ---------- Listings: Assembler der Registermaschine und Java ----------
\lstset{basicstyle=\ttfamily\small,keywordstyle=\bfseries,commentstyle=\color{gray!80!black},stringstyle=\color{black},showstringspaces=false,columns=flexible,keepspaces=true,tabsize=3,frame=single,framerule=0.4pt,rulecolor=\color{black},xleftmargin=2mm,framexleftmargin=1mm,aboveskip=1.5mm,belowskip=1.5mm,breaklines=true,escapeinside={(*}{*)},
 literate={ä}{{\"a}}1 {ö}{{\"o}}1 {ü}{{\"u}}1 {Ä}{{\"A}}1 {Ö}{{\"O}}1 {Ü}{{\"U}}1 {ß}{{\ss}}1 {–}{{--}}1 {…}{{\dots}}1 {≤}{{$\le$}}1 {≥}{{$\ge$}}1 {≠}{{$\ne$}}1}
\lstdefinelanguage{rm}{morekeywords={LOAD,LOADI,STORE,ADD,ADDI,SUB,SUBI,MUL,MULI,DIV,DIVI,MOD,MODI,CMP,CMPI,JMP,JEQ,JNE,JGT,JGE,JLT,JLE,HOLD},sensitive=true,morecomment=[l]{;}}
\lstnewenvironment{asm}[1][]{\lstset{language=rm,columns=fixed,numbers=left,firstnumber=0,numberstyle=\scriptsize\color{gray!70!black},numbersep=2.5mm,xleftmargin=6mm,framexleftmargin=5.5mm,#1}}{}
\lstnewenvironment{asmluecke}[1][]{\lstset{language=rm,columns=fixed,numbers=left,firstnumber=0,numberstyle=\scriptsize\color{gray!70!black},numbersep=2.5mm,xleftmargin=6mm,framexleftmargin=5.5mm,basicstyle=\linespread{1.5}\small\ttfamily,#1}}{}
\lstnewenvironment{java}[1][]{\lstset{language=Java,#1}}{}

% Sprungmarke in \lsgfeld-Lösungen (feste Breite, damit die Befehle untereinander stehen): \lb{kopf:}LOAD 100
\newcommand{\lb}[1]{\makebox[7ex][l]{#1}}

% ---------- Struktogramme (TikZ, Maße in mm, Zeilenhöhe 6 mm) ----------
%  \sgwenn[Breite]{Bedingung}{Wahr-Zweig}{Falsch-Zweig}
%  \sgschleife[Breite]{Anw. davor, …}{Schleifenkopf}{Rumpf-Anw., …}   (Listen mit Komma getrennt)
%  Farbe: aktuelle Textfarbe (in Lösungen z. B. {\lsgfarbe\sgwenn{…}})
\newcommand{\sgwenn}[4][60]{\begin{tikzpicture}[x=1mm,y=1mm,font=\small,line width=0.6pt,baseline=(current bounding box.north)]
  \draw (0,0) rectangle (#1,-15); \draw (0,-8.5) -- (#1,-8.5); \draw (#1/2,-8.5) -- (#1/2,-15);
  \draw[line width=0.4pt] (0,0) -- (#1/2,-8.5) -- (#1,0);
  \node at (#1/2,-2.8) {#2};
  \node[anchor=south west,font=\scriptsize,inner sep=0.6mm] at (0,-8.5) {wahr};
  \node[anchor=south east,font=\scriptsize,inner sep=0.6mm] at (#1,-8.5) {falsch};
  \node at (#1/4,-11.75) {#3}; \node at (3*#1/4,-11.75) {#4};
\end{tikzpicture}}
\newcommand{\sgschleife}[4][60]{\begin{tikzpicture}[x=1mm,y=1mm,font=\small,line width=0.6pt,baseline=(current bounding box.north)]
  \gdef\sgn{0}\gdef\sgm{0}
  \foreach \t [count=\i] in {#2}{\node[anchor=west] at (1.5,-6*\i+3) {\t}; \draw (0,-6*\i) -- (#1,-6*\i); \xdef\sgn{\i}}
  \node[anchor=west] at (1.5,-6*\sgn-3) {#3};
  \foreach \t [count=\i] in {#4}{\node[anchor=west] at (6.5,-6*\sgn-6*\i-3) {\t}; \draw (5,-6*\sgn-6*\i) -- (#1,-6*\sgn-6*\i); \xdef\sgm{\i}}
  \draw (0,0) rectangle (#1,{-6*(\sgn+1+\sgm)});
  \draw (5,{-6*(\sgn+1+\sgm)}) -- (5,{-6*(\sgn+1)});
\end{tikzpicture}}

% ---------- Farben wie auf der Website ----------
\definecolor{steuer}{HTML}{059669}\definecolor{rechen}{HTML}{EA580C}\definecolor{speicher}{HTML}{2563EB}
\definecolor{ea}{HTML}{7C3AED}\definecolor{busadr}{HTML}{D97706}\definecolor{busdaten}{HTML}{0891B2}\definecolor{bussteuer}{HTML}{BE123C}
\setlength{\unitlength}{1mm}

%  Übersetzen mit XeLaTeX (oder Tectonic):
%
%  Blatt:   xelatex ab-2-2-registermaschine.tex
%  Lösung:  xelatex ab-2-2-registermaschine-lsg.tex
%           (Datei enthält nur: \def\mitloesung{}\documentclass[11pt,a4paper]{article}
\input{ab-vorlage}
% Blatt:   xelatex ab-2-2-registermaschine.tex
% Lösung:  xelatex ab-2-2-registermaschine-lsg.tex
% Progression: erstes Programm vorhersagen und im Simulator Takt für Takt beobachten (Register, Busse)
%              -> Merkkästen Register und Befehlszyklus ①②③ -> Ablauf protokollieren
%              -> Statusregister N, Z, O -> selbst bauen: zwei Zellen tauschen (Hilfszelle)
\begin{document}
\abtitel{2.2}{Die Registermaschine}

Die Registermaschine ist ein vereinfachtes Modell eines Von-Neumann-Rechners. (x) bedeutet „Inhalt von Zelle x“.
Befehle mit I am Ende rechnen mit der Zahl im Befehl: \code{ADD 5} addiert (5), \code{ADDI 5} addiert die Zahl 5.

\begin{aufgabe}[Ein erstes Programm]{1}
\begin{minipage}[t]{0.36\linewidth}
\vspace{0pt}
\begin{asm}
LOAD 100
ADD 101
STORE 102
HOLD
\end{asm}
{\small Vorher: (100) = 7, (101) = 5}
\end{minipage}\hfill
\begin{minipage}[t]{0.6\linewidth}
\vspace{0pt}\lueckentext
\textbf{a)} Nach dem Programm: (102) = \lsg[2.25cm]{12} \quad Akkumulator: \lsg[2.25cm]{12}\\
\textbf{b)} Beobachte im Simulator den Befehl \code{ADD 101} Takt für Takt.
Was ändert sich in welcher Phase, was liegt auf den Bussen?
\end{minipage}\\[1mm]
{\renewcommand{\arraystretch}{1.05}\small
\begin{tabularx}{\linewidth}{|l|>{\raggedright\arraybackslash}X|p{2.2cm}|p{2.2cm}|p{2.6cm}|}\hline
\textbf{Phase} & \textbf{Diese Register ändern sich} & \textbf{Adressbus} & \textbf{Steuerbus} & \textbf{Datenbus}\\\hline
\nr{1} Holen\schreibzeile & \lsgtext{Befehlsregister (ADD 101), Befehlszähler (1 $\to$ 2)} & \lsgtext{1} & \lsgtext{lesen} & \lsgtext{ADD 101}\\\hline
\nr{2} Dekodieren\schreibzeile & \lsgtext{keines} & \lsgtext{–} & \lsgtext{–} & \lsgtext{–}\\\hline
\nr{3} Ausführen\schreibzeile & \lsgtext{Akkumulator (12), Statusregister} & \lsgtext{101} & \lsgtext{lesen} & \lsgtext{5}\\\hline
\end{tabularx}}
\end{aufgabe}

\begin{merke}[Register der Registermaschine]
{\renewcommand{\arraystretch}{1.05}
\begin{tabularx}{\linewidth}{|l|l|>{\raggedright\arraybackslash}X|}\hline
\textbf{Register} & \textbf{gehört zum} & \textbf{enthält}\\\hline
Befehlszähler (BZ)\schreibzeile & Steuerwerk & \lsgtext{die Adresse des nächsten Befehls}\\\hline
Befehlsregister (BR)\schreibzeile & Steuerwerk & \lsgtext{den Befehl, der gerade ausgeführt wird}\\\hline
Akkumulator (A)\schreibzeile & Rechenwerk & \lsgtext{einen Operanden und danach das Ergebnis jeder Rechnung}\\\hline
Statusregister (SR)\rule[-6.5mm]{0pt}{12mm} & Rechenwerk & \lsgtext{Eigenschaften des letzten Ergebnisses: N = 1 negativ, Z = 1 null, O = 1 Überlauf}\\\hline
\end{tabularx}}
\end{merke}

\begin{merke}[Der Befehlszyklus (wiederholt sich bis HOLD)]
\lueckentext
\nr{1} \textbf{Holen:} Der Befehl aus der Zelle, deren Adresse im \lsg[5.25cm]{Befehlszähler} steht, kommt über den Datenbus
ins \lsg[5.25cm]{Befehlsregister}. Der Befehlszähler wird \lsg[5.25cm]{um 1 erhöht}.\\
\nr{2} \textbf{Dekodieren:} Das Steuerwerk \lsg[5.25cm]{entschlüsselt} den Befehl.\\
\nr{3} \textbf{Ausführen:} Falls nötig wird ein Operand über die Busse geholt, das Rechenwerk rechnet, und das Ergebnis
kommt in den \lsg[5.25cm]{Akkumulator} (bei \code{STORE} in den Speicher).
\end{merke}

\newpage
\begin{aufgabe}[Ablauf protokollieren \normalfont(zuerst ohne Simulator)]{2}
\begin{minipage}[t]{0.3\linewidth}
\vspace{0pt}
\begin{asm}
LOAD 100
SUBI 6
STORE 101
LOADI 4
SUB 100
MULI 3
STORE 102
HOLD
\end{asm}
{\small Vorher: (100) = 6}
\end{minipage}\hfill
\begin{minipage}[t]{0.66\linewidth}
\vspace{0pt}
{\renewcommand{\arraystretch}{1}\small
\begin{tabularx}{\linewidth}{|c|l|c|c|c|X|}\hline
\textbf{Zelle} & \textbf{Befehl} & \textbf{A} & \textbf{N} & \textbf{Z} & \textbf{Speicher}\\\hline
0 & LOAD 100\schreibzeile & \lsgtext{6} & \lsgtext{0} & \lsgtext{0} & \\\hline
1 & SUBI 6\schreibzeile & \lsgtext{0} & \lsgtext{0} & \lsgtext{1} & \\\hline
2 & STORE 101\schreibzeile & \lsgtext{0} & \lsgtext{0} & \lsgtext{1} & \lsgtext{(101) = 0}\\\hline
3 & LOADI 4\schreibzeile & \lsgtext{4} & \lsgtext{0} & \lsgtext{0} & \\\hline
4 & SUB 100\schreibzeile & \lsgtext{−2} & \lsgtext{1} & \lsgtext{0} & \\\hline
5 & MULI 3\schreibzeile & \lsgtext{−6} & \lsgtext{1} & \lsgtext{0} & \\\hline
6 & STORE 102\schreibzeile & \lsgtext{−6} & \lsgtext{1} & \lsgtext{0} & \lsgtext{(102) = −6}\\\hline
\end{tabularx}}
\end{minipage}\\[1mm]
\lueckentext
Nach \code{STORE 101} ist Z immer noch 1. Warum? \lsgrest{STORE rechnet nicht: Das SR bleibt wie nach SUBI 6.}
\end{aufgabe}

\begin{aufgabe}[Das Statusregister \normalfont(Prüfen im Flaggen-Labor)]{3}
Ganze Zahlen haben 16 Bit: −32\,768 bis 32\,767. Was steht nach dem Befehl im Akkumulator, wie sind N, Z, O gesetzt?\\[1mm]
{\renewcommand{\arraystretch}{1}\small
\begin{tabularx}{\linewidth}{|c|c|>{\centering\arraybackslash}X|c|c|c|}\hline
\textbf{A vorher} & \textbf{Befehl} & \textbf{A danach} & \textbf{N} & \textbf{Z} & \textbf{O}\\\hline
5\schreibzeile & SUBI 5 & \lsgtext{0} & \lsgtext{0} & \lsgtext{1} & \lsgtext{0}\\\hline
5\schreibzeile & SUBI 8 & \lsgtext{−3} & \lsgtext{1} & \lsgtext{0} & \lsgtext{0}\\\hline
−3\schreibzeile & MULI −2 & \lsgtext{6} & \lsgtext{0} & \lsgtext{0} & \lsgtext{0}\\\hline
17\schreibzeile & MODI 5 & \lsgtext{2} & \lsgtext{0} & \lsgtext{0} & \lsgtext{0}\\\hline
32\,767\schreibzeile & ADDI 1 & \lsgtext{−32\,768 (Überlauf, falscher Wert)} & \lsgtext{1} & \lsgtext{0} & \lsgtext{1}\\\hline
\end{tabularx}}
\end{aufgabe}

\begin{aufgabe}[Selbst bauen: zwei Zellen tauschen \normalfont(dann im Simulator testen)]{4}
\begin{minipage}[t]{0.42\linewidth}
\vspace{0pt}\raggedright
Die Inhalte der Zellen 100 und 101 sollen vertauscht werden. Zelle 102 darfst du als Hilfszelle benutzen.\\[2mm]
Warum braucht man die Hilfszelle?
\lsglinien{3}{Der Akkumulator hält nur eine Zahl. Ohne Sicherung wäre der alte Inhalt von 100 beim ersten STORE verloren.}
\end{minipage}\hfill
\begin{minipage}[t]{0.54\linewidth}
\vspace{0pt}
\lsgfeld{7}{LOAD 100\\STORE 102 \ \ ; hilf = a\\LOAD 101\\STORE 100 \ \ ; a = b\\LOAD 102\\STORE 101 \ \ ; b = hilf\\HOLD}
\end{minipage}
\end{aufgabe}
\end{document}
)
%  Lösungseinträge belegen nie mehr Platz als die Lücke im Blatt.
%
%  Antwortplatz:
%   \luecke[l] / \lsg[l]{Text}      Lücke (Standard 6,75 cm, Fachbegriff 5,25 cm, Zahl 2,25 cm)
%   \restzeile / \lsgrest{Text}     Lücke bis Zeilenende      \lueckentext  Absatz mit Lücken
%   \linien{n} / \lsglinien{n}{T}   Schreiblinien 8 mm        \linienfeld{n} / \lsgfeld{n}{Code}
%   \schreibzeile                   Tabellenzelle 8 mm        \lsgtext{…}  nur in der Lösung
%   \kariert{n} / \lsgkariert{n}{TikZ}, \karobox{b}{h} / \lsgkarobox{b}{h}{TikZ}
%  Assembler:
%   asm           Listing, Zeilennummern = Adressen (ab 0); keine Leerzeilen, Marken in der Befehlszeile
%   asmluecke     dasselbe mit 1,5-fachem Zeilenabstand für Lücken (*\lsg[3.5cm]{JGE ende}*)
%   \nr{1}        Schrittnummer ① wie auf der Website und im Simulator
%  Struktogramme: \sgwenn, \sgschleife (siehe unten)
% =====================================================================
\usepackage[a4paper,left=16mm,right=16mm,top=23mm,bottom=13mm,headheight=12mm,headsep=3mm,footskip=6mm]{geometry}
\usepackage{iftex}
\ifPDFTeX
  \usepackage[T1]{fontenc}
  \usepackage[utf8]{inputenc}
  \usepackage{helvet}
  \usepackage[scaled=0.86]{beramono}
  \renewcommand{\familydefault}{\sfdefault}
\else
  \usepackage{fontspec}
  \setmainfont{texgyreheros}[Extension=.otf,UprightFont=*-regular,BoldFont=*-bold,ItalicFont=*-italic,BoldItalicFont=*-bolditalic]
  \setmonofont{DejaVuSansMono}[Extension=.ttf,UprightFont=*,BoldFont=*-Bold,ItalicFont=*-Oblique,BoldItalicFont=*-BoldOblique,Scale=0.86]
\fi
\usepackage[ngerman,shorthands=off]{babel}
\usepackage{amssymb}
\usepackage{xcolor,graphicx,tabularx,array,enumitem,fancyhdr,tikz,multicol,calc,etoolbox,pifont,xspace,listings}
\usepackage[most]{tcolorbox}
\usetikzlibrary{arrows.meta,positioning,calc}
\usepackage[hidelinks]{hyperref}
\setlength{\parindent}{0pt}
\setlength{\parskip}{3pt}
\setlist{nosep,leftmargin=*}

% ---------- Lösungsschalter ----------
\newif\ifloesung
\ifdefined\mitloesung\loesungtrue\fi
\newcommand{\lsgfarbe}{\color{red!75!black}}

% ---------- Kopf- und Fußzeile ----------
\newcommand{\lehrkraft}{Hau}
\newcommand{\themenbereich}{2. Funktionsweise eines Rechners}
\newcommand{\abnummer}{}
\pagestyle{fancy}\fancyhf{}
\renewcommand{\headrulewidth}{0.4pt}
\fancyhead[L]{\small Klasse 13\\Datum:}
\fancyhead[C]{\small Themenbereich:\\\themenbereich}
\fancyhead[R]{\small \lehrkraft\ifloesung\\{\lsgfarbe\bfseries Lösung}\fi}
\fancyfoot[C]{\scriptsize edu-mrh.de\,·\,Informatik 13\,·\,Blatt \abnummer\ifloesung\ (Lösung)\fi\,·\,Seite \thepage}
\newcommand{\abtitel}[2]{\renewcommand{\abnummer}{#1}\begin{center}\Large\bfseries #1\quad\underline{#2}\end{center}\vspace{-2mm}}

% ---------- Boxen ----------
\newenvironment{aufgabe}[2][]{\begin{tcolorbox}[enhanced,breakable,colback=white,colframe=black,boxrule=0.7pt,arc=3mm,left=2mm,right=2mm,top=1mm,bottom=1.2mm,before skip=2mm,after skip=2mm]\textbf{Aufgabe #2}\ifstrempty{#1}{}{ \textit{#1}}\par\vspace{0.3mm}}{\end{tcolorbox}}
\newtcolorbox{merke}[1][Merke]{enhanced,breakable,colback=green!5,colframe=green!40!black,boxrule=0.8pt,arc=2mm,left=2mm,right=2mm,top=1mm,bottom=1.2mm,before skip=2mm,after skip=2mm,title={#1},fonttitle=\bfseries,coltitle=white,toptitle=0.3mm,bottomtitle=0.3mm}
\newtcolorbox{hinweis}{enhanced,breakable,colback=yellow!12,colframe=orange!70!black,boxrule=0.6pt,arc=2mm,left=2mm,right=2mm,top=1mm,bottom=1mm,before skip=2mm,after skip=2mm}
\newcommand{\web}[1]{{\small\color{gray}\textit{Website: #1}}}
\newcommand{\code}[1]{\texttt{#1}}

% ---------- Antwortplatz (wie informatik/10/java, Lücken x 1,5) ----------
\newlength{\zeilenabstand}\setlength{\zeilenabstand}{8mm}
\newlength{\karo}\setlength{\karo}{5mm}
\tikzset{karo/.style={draw=gray!50,line width=0.3pt},schreiblinie/.style={draw=black!65,line width=0.35pt}}
\newlength{\lsgbreite}
\newcommand{\lsg}[2][6.75cm]{\leavevmode\ifloesung\settowidth{\lsgbreite}{\,#2\,}%
  \ifdim\lsgbreite>#1\relax\rlap{\raisebox{0.8pt}{\resizebox{#1}{!}{\lsgfarbe\,#2\,}}}%
  \else\rlap{\raisebox{0.8pt}{\lsgfarbe\,#2}}\fi\fi\rule[-0.6ex]{#1}{0.4pt}\xspace}
\newcommand{\luecke}[1][6.75cm]{\lsg[#1]{}}
\newcommand{\lsgrest}[1]{\leavevmode\unskip\hspace{1.5mm}\ifloesung\rlap{\raisebox{0.8pt}{\lsgfarbe\,#1}}\fi\leaders\hbox{\rule[-0.6ex]{0.5pt}{0.4pt}}\hfill\null\par}
\newcommand{\restzeile}{\lsgrest{}}
\newcommand{\lsgtext}[1]{\ifloesung\leavevmode{\lsgfarbe#1}\fi}
\newcommand{\lsgoder}[2]{\ifloesung\leavevmode#1\else#2\fi}
\newcommand{\lueckentext}{\linespread{1.5}\selectfont\raggedright}
\newcommand{\schreibzeile}{\rule[-2.8mm]{0pt}{8mm}}
\newcommand{\nr}[1]{\ding{\numexpr171+#1\relax}}
\newcommand{\lsglinien}[2]{\par\vspace{0.5mm}\noindent\begin{tikzpicture}
  \path (0,0) (\linewidth-0.4pt,-#1\zeilenabstand-1.5mm);
  \foreach \i in {1,...,#1}{\draw[schreiblinie] (0.2pt,-\i\zeilenabstand) -- (\linewidth-0.2pt,-\i\zeilenabstand);}
  \ifloesung\node[overlay,anchor=north west,inner sep=0pt] at (1mm,-\zeilenabstand+1pt+5mm) {\parbox[t]{\linewidth-2mm}{\lsgfarbe\raggedright\setlength{\baselineskip}{\zeilenabstand}\rule{0pt}{5mm}#2}};\fi
  \end{tikzpicture}\par}
\newcommand{\linien}[1]{\lsglinien{#1}{}}
\newcommand{\lsgfeld}[2]{\noindent\begin{tikzpicture}
  \draw[line width=0.4pt] (0.2pt,0) rectangle (\linewidth-0.2pt,-#1\zeilenabstand-3mm);
  \foreach \i in {1,...,#1}{\draw[schreiblinie] (2mm,-\i\zeilenabstand) -- (\linewidth-2mm,-\i\zeilenabstand);}
  \ifloesung\node[overlay,anchor=north west,inner sep=0pt] at (2.5mm,-\zeilenabstand+1pt+5mm) {\parbox[t]{\linewidth-5mm}{\lsgfarbe\ttfamily\small\setlength{\baselineskip}{\zeilenabstand}\rule{0pt}{5mm}#2}};\fi
  \end{tikzpicture}}
\newcommand{\linienfeld}[1]{\lsgfeld{#1}{}}
\newcommand{\karoraster}[3]{\begin{tikzpicture}\draw[karo] (0,0) grid[step=\karo] (#1\karo,#2\karo);
  \ifloesung\begin{scope}[overlay,shift={(0,#2\karo)},red!75!black,text=red!75!black,line width=0.8pt]#3\end{scope}\fi\end{tikzpicture}}
\newcommand{\lsgkariert}[2]{\par\vspace{1.5mm}\noindent\pgfmathtruncatemacro{\karoX}{floor(\linewidth/\karo)}%
  \makebox[\linewidth]{\karoraster{\karoX}{#1}{#2}}\par\vspace{1mm}}
\newcommand{\kariert}[1]{\lsgkariert{#1}{}}
\newcommand{\lsgkarobox}[3]{\pgfmathtruncatemacro{\karoX}{round((#1)/\karo)}\pgfmathtruncatemacro{\karoY}{round((#2)/\karo)}\karoraster{\karoX}{\karoY}{#3}}
\newcommand{\karobox}[2]{\lsgkarobox{#1}{#2}{}}
% Ankreuzfeld: \kreuz{1} = angekreuzt (nur in der Lösung)
\newcommand{\kreuz}[1]{\leavevmode\ifloesung\ifnum#1=1 {\lsgfarbe$\boxtimes$}\else$\square$\fi\else$\square$\fi}

% ---------- Listings: Assembler der Registermaschine und Java ----------
\lstset{basicstyle=\ttfamily\small,keywordstyle=\bfseries,commentstyle=\color{gray!80!black},stringstyle=\color{black},showstringspaces=false,columns=flexible,keepspaces=true,tabsize=3,frame=single,framerule=0.4pt,rulecolor=\color{black},xleftmargin=2mm,framexleftmargin=1mm,aboveskip=1.5mm,belowskip=1.5mm,breaklines=true,escapeinside={(*}{*)},
 literate={ä}{{\"a}}1 {ö}{{\"o}}1 {ü}{{\"u}}1 {Ä}{{\"A}}1 {Ö}{{\"O}}1 {Ü}{{\"U}}1 {ß}{{\ss}}1 {–}{{--}}1 {…}{{\dots}}1 {≤}{{$\le$}}1 {≥}{{$\ge$}}1 {≠}{{$\ne$}}1}
\lstdefinelanguage{rm}{morekeywords={LOAD,LOADI,STORE,ADD,ADDI,SUB,SUBI,MUL,MULI,DIV,DIVI,MOD,MODI,CMP,CMPI,JMP,JEQ,JNE,JGT,JGE,JLT,JLE,HOLD},sensitive=true,morecomment=[l]{;}}
\lstnewenvironment{asm}[1][]{\lstset{language=rm,columns=fixed,numbers=left,firstnumber=0,numberstyle=\scriptsize\color{gray!70!black},numbersep=2.5mm,xleftmargin=6mm,framexleftmargin=5.5mm,#1}}{}
\lstnewenvironment{asmluecke}[1][]{\lstset{language=rm,columns=fixed,numbers=left,firstnumber=0,numberstyle=\scriptsize\color{gray!70!black},numbersep=2.5mm,xleftmargin=6mm,framexleftmargin=5.5mm,basicstyle=\linespread{1.5}\small\ttfamily,#1}}{}
\lstnewenvironment{java}[1][]{\lstset{language=Java,#1}}{}

% Sprungmarke in \lsgfeld-Lösungen (feste Breite, damit die Befehle untereinander stehen): \lb{kopf:}LOAD 100
\newcommand{\lb}[1]{\makebox[7ex][l]{#1}}

% ---------- Struktogramme (TikZ, Maße in mm, Zeilenhöhe 6 mm) ----------
%  \sgwenn[Breite]{Bedingung}{Wahr-Zweig}{Falsch-Zweig}
%  \sgschleife[Breite]{Anw. davor, …}{Schleifenkopf}{Rumpf-Anw., …}   (Listen mit Komma getrennt)
%  Farbe: aktuelle Textfarbe (in Lösungen z. B. {\lsgfarbe\sgwenn{…}})
\newcommand{\sgwenn}[4][60]{\begin{tikzpicture}[x=1mm,y=1mm,font=\small,line width=0.6pt,baseline=(current bounding box.north)]
  \draw (0,0) rectangle (#1,-15); \draw (0,-8.5) -- (#1,-8.5); \draw (#1/2,-8.5) -- (#1/2,-15);
  \draw[line width=0.4pt] (0,0) -- (#1/2,-8.5) -- (#1,0);
  \node at (#1/2,-2.8) {#2};
  \node[anchor=south west,font=\scriptsize,inner sep=0.6mm] at (0,-8.5) {wahr};
  \node[anchor=south east,font=\scriptsize,inner sep=0.6mm] at (#1,-8.5) {falsch};
  \node at (#1/4,-11.75) {#3}; \node at (3*#1/4,-11.75) {#4};
\end{tikzpicture}}
\newcommand{\sgschleife}[4][60]{\begin{tikzpicture}[x=1mm,y=1mm,font=\small,line width=0.6pt,baseline=(current bounding box.north)]
  \gdef\sgn{0}\gdef\sgm{0}
  \foreach \t [count=\i] in {#2}{\node[anchor=west] at (1.5,-6*\i+3) {\t}; \draw (0,-6*\i) -- (#1,-6*\i); \xdef\sgn{\i}}
  \node[anchor=west] at (1.5,-6*\sgn-3) {#3};
  \foreach \t [count=\i] in {#4}{\node[anchor=west] at (6.5,-6*\sgn-6*\i-3) {\t}; \draw (5,-6*\sgn-6*\i) -- (#1,-6*\sgn-6*\i); \xdef\sgm{\i}}
  \draw (0,0) rectangle (#1,{-6*(\sgn+1+\sgm)});
  \draw (5,{-6*(\sgn+1+\sgm)}) -- (5,{-6*(\sgn+1)});
\end{tikzpicture}}

% ---------- Farben wie auf der Website ----------
\definecolor{steuer}{HTML}{059669}\definecolor{rechen}{HTML}{EA580C}\definecolor{speicher}{HTML}{2563EB}
\definecolor{ea}{HTML}{7C3AED}\definecolor{busadr}{HTML}{D97706}\definecolor{busdaten}{HTML}{0891B2}\definecolor{bussteuer}{HTML}{BE123C}
\setlength{\unitlength}{1mm}

%  Übersetzen mit XeLaTeX (oder Tectonic):
%
%  Blatt:   xelatex ab-2-2-registermaschine.tex
%  Lösung:  xelatex ab-2-2-registermaschine-lsg.tex
%           (Datei enthält nur: \def\mitloesung{}\documentclass[11pt,a4paper]{article}
% =====================================================================
%  Gemeinsame Vorlage der Arbeitsblätter – 2. Funktionsweise eines Rechners (Q13)
%  Kopf, Fuß und Boxen wie informatik/13/formale-sprachen/arbeitsblaetter/ab-vorlage.tex,
%  Antwortplatz-Makros wie informatik/10/java/arbeitsblaetter/ab-vorlage.tex
%  (Linien 8 mm, Karo 5 mm aus ganzen Kästchen, Lücken seit 2026-10 x 1,5).
%  Einbinden in jedem Blatt direkt nach \documentclass: \input{ab-vorlage}
%  Übersetzen mit XeLaTeX (oder Tectonic):
%
%  Blatt:   xelatex ab-2-2-registermaschine.tex
%  Lösung:  xelatex ab-2-2-registermaschine-lsg.tex
%           (Datei enthält nur: \def\mitloesung{}\input{ab-2-2-registermaschine})
%  Lösungseinträge belegen nie mehr Platz als die Lücke im Blatt.
%
%  Antwortplatz:
%   \luecke[l] / \lsg[l]{Text}      Lücke (Standard 6,75 cm, Fachbegriff 5,25 cm, Zahl 2,25 cm)
%   \restzeile / \lsgrest{Text}     Lücke bis Zeilenende      \lueckentext  Absatz mit Lücken
%   \linien{n} / \lsglinien{n}{T}   Schreiblinien 8 mm        \linienfeld{n} / \lsgfeld{n}{Code}
%   \schreibzeile                   Tabellenzelle 8 mm        \lsgtext{…}  nur in der Lösung
%   \kariert{n} / \lsgkariert{n}{TikZ}, \karobox{b}{h} / \lsgkarobox{b}{h}{TikZ}
%  Assembler:
%   asm           Listing, Zeilennummern = Adressen (ab 0); keine Leerzeilen, Marken in der Befehlszeile
%   asmluecke     dasselbe mit 1,5-fachem Zeilenabstand für Lücken (*\lsg[3.5cm]{JGE ende}*)
%   \nr{1}        Schrittnummer ① wie auf der Website und im Simulator
%  Struktogramme: \sgwenn, \sgschleife (siehe unten)
% =====================================================================
\usepackage[a4paper,left=16mm,right=16mm,top=23mm,bottom=13mm,headheight=12mm,headsep=3mm,footskip=6mm]{geometry}
\usepackage{iftex}
\ifPDFTeX
  \usepackage[T1]{fontenc}
  \usepackage[utf8]{inputenc}
  \usepackage{helvet}
  \usepackage[scaled=0.86]{beramono}
  \renewcommand{\familydefault}{\sfdefault}
\else
  \usepackage{fontspec}
  \setmainfont{texgyreheros}[Extension=.otf,UprightFont=*-regular,BoldFont=*-bold,ItalicFont=*-italic,BoldItalicFont=*-bolditalic]
  \setmonofont{DejaVuSansMono}[Extension=.ttf,UprightFont=*,BoldFont=*-Bold,ItalicFont=*-Oblique,BoldItalicFont=*-BoldOblique,Scale=0.86]
\fi
\usepackage[ngerman,shorthands=off]{babel}
\usepackage{amssymb}
\usepackage{xcolor,graphicx,tabularx,array,enumitem,fancyhdr,tikz,multicol,calc,etoolbox,pifont,xspace,listings}
\usepackage[most]{tcolorbox}
\usetikzlibrary{arrows.meta,positioning,calc}
\usepackage[hidelinks]{hyperref}
\setlength{\parindent}{0pt}
\setlength{\parskip}{3pt}
\setlist{nosep,leftmargin=*}

% ---------- Lösungsschalter ----------
\newif\ifloesung
\ifdefined\mitloesung\loesungtrue\fi
\newcommand{\lsgfarbe}{\color{red!75!black}}

% ---------- Kopf- und Fußzeile ----------
\newcommand{\lehrkraft}{Hau}
\newcommand{\themenbereich}{2. Funktionsweise eines Rechners}
\newcommand{\abnummer}{}
\pagestyle{fancy}\fancyhf{}
\renewcommand{\headrulewidth}{0.4pt}
\fancyhead[L]{\small Klasse 13\\Datum:}
\fancyhead[C]{\small Themenbereich:\\\themenbereich}
\fancyhead[R]{\small \lehrkraft\ifloesung\\{\lsgfarbe\bfseries Lösung}\fi}
\fancyfoot[C]{\scriptsize edu-mrh.de\,·\,Informatik 13\,·\,Blatt \abnummer\ifloesung\ (Lösung)\fi\,·\,Seite \thepage}
\newcommand{\abtitel}[2]{\renewcommand{\abnummer}{#1}\begin{center}\Large\bfseries #1\quad\underline{#2}\end{center}\vspace{-2mm}}

% ---------- Boxen ----------
\newenvironment{aufgabe}[2][]{\begin{tcolorbox}[enhanced,breakable,colback=white,colframe=black,boxrule=0.7pt,arc=3mm,left=2mm,right=2mm,top=1mm,bottom=1.2mm,before skip=2mm,after skip=2mm]\textbf{Aufgabe #2}\ifstrempty{#1}{}{ \textit{#1}}\par\vspace{0.3mm}}{\end{tcolorbox}}
\newtcolorbox{merke}[1][Merke]{enhanced,breakable,colback=green!5,colframe=green!40!black,boxrule=0.8pt,arc=2mm,left=2mm,right=2mm,top=1mm,bottom=1.2mm,before skip=2mm,after skip=2mm,title={#1},fonttitle=\bfseries,coltitle=white,toptitle=0.3mm,bottomtitle=0.3mm}
\newtcolorbox{hinweis}{enhanced,breakable,colback=yellow!12,colframe=orange!70!black,boxrule=0.6pt,arc=2mm,left=2mm,right=2mm,top=1mm,bottom=1mm,before skip=2mm,after skip=2mm}
\newcommand{\web}[1]{{\small\color{gray}\textit{Website: #1}}}
\newcommand{\code}[1]{\texttt{#1}}

% ---------- Antwortplatz (wie informatik/10/java, Lücken x 1,5) ----------
\newlength{\zeilenabstand}\setlength{\zeilenabstand}{8mm}
\newlength{\karo}\setlength{\karo}{5mm}
\tikzset{karo/.style={draw=gray!50,line width=0.3pt},schreiblinie/.style={draw=black!65,line width=0.35pt}}
\newlength{\lsgbreite}
\newcommand{\lsg}[2][6.75cm]{\leavevmode\ifloesung\settowidth{\lsgbreite}{\,#2\,}%
  \ifdim\lsgbreite>#1\relax\rlap{\raisebox{0.8pt}{\resizebox{#1}{!}{\lsgfarbe\,#2\,}}}%
  \else\rlap{\raisebox{0.8pt}{\lsgfarbe\,#2}}\fi\fi\rule[-0.6ex]{#1}{0.4pt}\xspace}
\newcommand{\luecke}[1][6.75cm]{\lsg[#1]{}}
\newcommand{\lsgrest}[1]{\leavevmode\unskip\hspace{1.5mm}\ifloesung\rlap{\raisebox{0.8pt}{\lsgfarbe\,#1}}\fi\leaders\hbox{\rule[-0.6ex]{0.5pt}{0.4pt}}\hfill\null\par}
\newcommand{\restzeile}{\lsgrest{}}
\newcommand{\lsgtext}[1]{\ifloesung\leavevmode{\lsgfarbe#1}\fi}
\newcommand{\lsgoder}[2]{\ifloesung\leavevmode#1\else#2\fi}
\newcommand{\lueckentext}{\linespread{1.5}\selectfont\raggedright}
\newcommand{\schreibzeile}{\rule[-2.8mm]{0pt}{8mm}}
\newcommand{\nr}[1]{\ding{\numexpr171+#1\relax}}
\newcommand{\lsglinien}[2]{\par\vspace{0.5mm}\noindent\begin{tikzpicture}
  \path (0,0) (\linewidth-0.4pt,-#1\zeilenabstand-1.5mm);
  \foreach \i in {1,...,#1}{\draw[schreiblinie] (0.2pt,-\i\zeilenabstand) -- (\linewidth-0.2pt,-\i\zeilenabstand);}
  \ifloesung\node[overlay,anchor=north west,inner sep=0pt] at (1mm,-\zeilenabstand+1pt+5mm) {\parbox[t]{\linewidth-2mm}{\lsgfarbe\raggedright\setlength{\baselineskip}{\zeilenabstand}\rule{0pt}{5mm}#2}};\fi
  \end{tikzpicture}\par}
\newcommand{\linien}[1]{\lsglinien{#1}{}}
\newcommand{\lsgfeld}[2]{\noindent\begin{tikzpicture}
  \draw[line width=0.4pt] (0.2pt,0) rectangle (\linewidth-0.2pt,-#1\zeilenabstand-3mm);
  \foreach \i in {1,...,#1}{\draw[schreiblinie] (2mm,-\i\zeilenabstand) -- (\linewidth-2mm,-\i\zeilenabstand);}
  \ifloesung\node[overlay,anchor=north west,inner sep=0pt] at (2.5mm,-\zeilenabstand+1pt+5mm) {\parbox[t]{\linewidth-5mm}{\lsgfarbe\ttfamily\small\setlength{\baselineskip}{\zeilenabstand}\rule{0pt}{5mm}#2}};\fi
  \end{tikzpicture}}
\newcommand{\linienfeld}[1]{\lsgfeld{#1}{}}
\newcommand{\karoraster}[3]{\begin{tikzpicture}\draw[karo] (0,0) grid[step=\karo] (#1\karo,#2\karo);
  \ifloesung\begin{scope}[overlay,shift={(0,#2\karo)},red!75!black,text=red!75!black,line width=0.8pt]#3\end{scope}\fi\end{tikzpicture}}
\newcommand{\lsgkariert}[2]{\par\vspace{1.5mm}\noindent\pgfmathtruncatemacro{\karoX}{floor(\linewidth/\karo)}%
  \makebox[\linewidth]{\karoraster{\karoX}{#1}{#2}}\par\vspace{1mm}}
\newcommand{\kariert}[1]{\lsgkariert{#1}{}}
\newcommand{\lsgkarobox}[3]{\pgfmathtruncatemacro{\karoX}{round((#1)/\karo)}\pgfmathtruncatemacro{\karoY}{round((#2)/\karo)}\karoraster{\karoX}{\karoY}{#3}}
\newcommand{\karobox}[2]{\lsgkarobox{#1}{#2}{}}
% Ankreuzfeld: \kreuz{1} = angekreuzt (nur in der Lösung)
\newcommand{\kreuz}[1]{\leavevmode\ifloesung\ifnum#1=1 {\lsgfarbe$\boxtimes$}\else$\square$\fi\else$\square$\fi}

% ---------- Listings: Assembler der Registermaschine und Java ----------
\lstset{basicstyle=\ttfamily\small,keywordstyle=\bfseries,commentstyle=\color{gray!80!black},stringstyle=\color{black},showstringspaces=false,columns=flexible,keepspaces=true,tabsize=3,frame=single,framerule=0.4pt,rulecolor=\color{black},xleftmargin=2mm,framexleftmargin=1mm,aboveskip=1.5mm,belowskip=1.5mm,breaklines=true,escapeinside={(*}{*)},
 literate={ä}{{\"a}}1 {ö}{{\"o}}1 {ü}{{\"u}}1 {Ä}{{\"A}}1 {Ö}{{\"O}}1 {Ü}{{\"U}}1 {ß}{{\ss}}1 {–}{{--}}1 {…}{{\dots}}1 {≤}{{$\le$}}1 {≥}{{$\ge$}}1 {≠}{{$\ne$}}1}
\lstdefinelanguage{rm}{morekeywords={LOAD,LOADI,STORE,ADD,ADDI,SUB,SUBI,MUL,MULI,DIV,DIVI,MOD,MODI,CMP,CMPI,JMP,JEQ,JNE,JGT,JGE,JLT,JLE,HOLD},sensitive=true,morecomment=[l]{;}}
\lstnewenvironment{asm}[1][]{\lstset{language=rm,columns=fixed,numbers=left,firstnumber=0,numberstyle=\scriptsize\color{gray!70!black},numbersep=2.5mm,xleftmargin=6mm,framexleftmargin=5.5mm,#1}}{}
\lstnewenvironment{asmluecke}[1][]{\lstset{language=rm,columns=fixed,numbers=left,firstnumber=0,numberstyle=\scriptsize\color{gray!70!black},numbersep=2.5mm,xleftmargin=6mm,framexleftmargin=5.5mm,basicstyle=\linespread{1.5}\small\ttfamily,#1}}{}
\lstnewenvironment{java}[1][]{\lstset{language=Java,#1}}{}

% Sprungmarke in \lsgfeld-Lösungen (feste Breite, damit die Befehle untereinander stehen): \lb{kopf:}LOAD 100
\newcommand{\lb}[1]{\makebox[7ex][l]{#1}}

% ---------- Struktogramme (TikZ, Maße in mm, Zeilenhöhe 6 mm) ----------
%  \sgwenn[Breite]{Bedingung}{Wahr-Zweig}{Falsch-Zweig}
%  \sgschleife[Breite]{Anw. davor, …}{Schleifenkopf}{Rumpf-Anw., …}   (Listen mit Komma getrennt)
%  Farbe: aktuelle Textfarbe (in Lösungen z. B. {\lsgfarbe\sgwenn{…}})
\newcommand{\sgwenn}[4][60]{\begin{tikzpicture}[x=1mm,y=1mm,font=\small,line width=0.6pt,baseline=(current bounding box.north)]
  \draw (0,0) rectangle (#1,-15); \draw (0,-8.5) -- (#1,-8.5); \draw (#1/2,-8.5) -- (#1/2,-15);
  \draw[line width=0.4pt] (0,0) -- (#1/2,-8.5) -- (#1,0);
  \node at (#1/2,-2.8) {#2};
  \node[anchor=south west,font=\scriptsize,inner sep=0.6mm] at (0,-8.5) {wahr};
  \node[anchor=south east,font=\scriptsize,inner sep=0.6mm] at (#1,-8.5) {falsch};
  \node at (#1/4,-11.75) {#3}; \node at (3*#1/4,-11.75) {#4};
\end{tikzpicture}}
\newcommand{\sgschleife}[4][60]{\begin{tikzpicture}[x=1mm,y=1mm,font=\small,line width=0.6pt,baseline=(current bounding box.north)]
  \gdef\sgn{0}\gdef\sgm{0}
  \foreach \t [count=\i] in {#2}{\node[anchor=west] at (1.5,-6*\i+3) {\t}; \draw (0,-6*\i) -- (#1,-6*\i); \xdef\sgn{\i}}
  \node[anchor=west] at (1.5,-6*\sgn-3) {#3};
  \foreach \t [count=\i] in {#4}{\node[anchor=west] at (6.5,-6*\sgn-6*\i-3) {\t}; \draw (5,-6*\sgn-6*\i) -- (#1,-6*\sgn-6*\i); \xdef\sgm{\i}}
  \draw (0,0) rectangle (#1,{-6*(\sgn+1+\sgm)});
  \draw (5,{-6*(\sgn+1+\sgm)}) -- (5,{-6*(\sgn+1)});
\end{tikzpicture}}

% ---------- Farben wie auf der Website ----------
\definecolor{steuer}{HTML}{059669}\definecolor{rechen}{HTML}{EA580C}\definecolor{speicher}{HTML}{2563EB}
\definecolor{ea}{HTML}{7C3AED}\definecolor{busadr}{HTML}{D97706}\definecolor{busdaten}{HTML}{0891B2}\definecolor{bussteuer}{HTML}{BE123C}
\setlength{\unitlength}{1mm}

% Blatt:   xelatex ab-2-2-registermaschine.tex
% Lösung:  xelatex ab-2-2-registermaschine-lsg.tex
% Progression: erstes Programm vorhersagen und im Simulator Takt für Takt beobachten (Register, Busse)
%              -> Merkkästen Register und Befehlszyklus ①②③ -> Ablauf protokollieren
%              -> Statusregister N, Z, O -> selbst bauen: zwei Zellen tauschen (Hilfszelle)
\begin{document}
\abtitel{2.2}{Die Registermaschine}

Die Registermaschine ist ein vereinfachtes Modell eines Von-Neumann-Rechners. (x) bedeutet „Inhalt von Zelle x“.
Befehle mit I am Ende rechnen mit der Zahl im Befehl: \code{ADD 5} addiert (5), \code{ADDI 5} addiert die Zahl 5.

\begin{aufgabe}[Ein erstes Programm]{1}
\begin{minipage}[t]{0.36\linewidth}
\vspace{0pt}
\begin{asm}
LOAD 100
ADD 101
STORE 102
HOLD
\end{asm}
{\small Vorher: (100) = 7, (101) = 5}
\end{minipage}\hfill
\begin{minipage}[t]{0.6\linewidth}
\vspace{0pt}\lueckentext
\textbf{a)} Nach dem Programm: (102) = \lsg[2.25cm]{12} \quad Akkumulator: \lsg[2.25cm]{12}\\
\textbf{b)} Beobachte im Simulator den Befehl \code{ADD 101} Takt für Takt.
Was ändert sich in welcher Phase, was liegt auf den Bussen?
\end{minipage}\\[1mm]
{\renewcommand{\arraystretch}{1.05}\small
\begin{tabularx}{\linewidth}{|l|>{\raggedright\arraybackslash}X|p{2.2cm}|p{2.2cm}|p{2.6cm}|}\hline
\textbf{Phase} & \textbf{Diese Register ändern sich} & \textbf{Adressbus} & \textbf{Steuerbus} & \textbf{Datenbus}\\\hline
\nr{1} Holen\schreibzeile & \lsgtext{Befehlsregister (ADD 101), Befehlszähler (1 $\to$ 2)} & \lsgtext{1} & \lsgtext{lesen} & \lsgtext{ADD 101}\\\hline
\nr{2} Dekodieren\schreibzeile & \lsgtext{keines} & \lsgtext{–} & \lsgtext{–} & \lsgtext{–}\\\hline
\nr{3} Ausführen\schreibzeile & \lsgtext{Akkumulator (12), Statusregister} & \lsgtext{101} & \lsgtext{lesen} & \lsgtext{5}\\\hline
\end{tabularx}}
\end{aufgabe}

\begin{merke}[Register der Registermaschine]
{\renewcommand{\arraystretch}{1.05}
\begin{tabularx}{\linewidth}{|l|l|>{\raggedright\arraybackslash}X|}\hline
\textbf{Register} & \textbf{gehört zum} & \textbf{enthält}\\\hline
Befehlszähler (BZ)\schreibzeile & Steuerwerk & \lsgtext{die Adresse des nächsten Befehls}\\\hline
Befehlsregister (BR)\schreibzeile & Steuerwerk & \lsgtext{den Befehl, der gerade ausgeführt wird}\\\hline
Akkumulator (A)\schreibzeile & Rechenwerk & \lsgtext{einen Operanden und danach das Ergebnis jeder Rechnung}\\\hline
Statusregister (SR)\rule[-6.5mm]{0pt}{12mm} & Rechenwerk & \lsgtext{Eigenschaften des letzten Ergebnisses: N = 1 negativ, Z = 1 null, O = 1 Überlauf}\\\hline
\end{tabularx}}
\end{merke}

\begin{merke}[Der Befehlszyklus (wiederholt sich bis HOLD)]
\lueckentext
\nr{1} \textbf{Holen:} Der Befehl aus der Zelle, deren Adresse im \lsg[5.25cm]{Befehlszähler} steht, kommt über den Datenbus
ins \lsg[5.25cm]{Befehlsregister}. Der Befehlszähler wird \lsg[5.25cm]{um 1 erhöht}.\\
\nr{2} \textbf{Dekodieren:} Das Steuerwerk \lsg[5.25cm]{entschlüsselt} den Befehl.\\
\nr{3} \textbf{Ausführen:} Falls nötig wird ein Operand über die Busse geholt, das Rechenwerk rechnet, und das Ergebnis
kommt in den \lsg[5.25cm]{Akkumulator} (bei \code{STORE} in den Speicher).
\end{merke}

\newpage
\begin{aufgabe}[Ablauf protokollieren \normalfont(zuerst ohne Simulator)]{2}
\begin{minipage}[t]{0.3\linewidth}
\vspace{0pt}
\begin{asm}
LOAD 100
SUBI 6
STORE 101
LOADI 4
SUB 100
MULI 3
STORE 102
HOLD
\end{asm}
{\small Vorher: (100) = 6}
\end{minipage}\hfill
\begin{minipage}[t]{0.66\linewidth}
\vspace{0pt}
{\renewcommand{\arraystretch}{1}\small
\begin{tabularx}{\linewidth}{|c|l|c|c|c|X|}\hline
\textbf{Zelle} & \textbf{Befehl} & \textbf{A} & \textbf{N} & \textbf{Z} & \textbf{Speicher}\\\hline
0 & LOAD 100\schreibzeile & \lsgtext{6} & \lsgtext{0} & \lsgtext{0} & \\\hline
1 & SUBI 6\schreibzeile & \lsgtext{0} & \lsgtext{0} & \lsgtext{1} & \\\hline
2 & STORE 101\schreibzeile & \lsgtext{0} & \lsgtext{0} & \lsgtext{1} & \lsgtext{(101) = 0}\\\hline
3 & LOADI 4\schreibzeile & \lsgtext{4} & \lsgtext{0} & \lsgtext{0} & \\\hline
4 & SUB 100\schreibzeile & \lsgtext{−2} & \lsgtext{1} & \lsgtext{0} & \\\hline
5 & MULI 3\schreibzeile & \lsgtext{−6} & \lsgtext{1} & \lsgtext{0} & \\\hline
6 & STORE 102\schreibzeile & \lsgtext{−6} & \lsgtext{1} & \lsgtext{0} & \lsgtext{(102) = −6}\\\hline
\end{tabularx}}
\end{minipage}\\[1mm]
\lueckentext
Nach \code{STORE 101} ist Z immer noch 1. Warum? \lsgrest{STORE rechnet nicht: Das SR bleibt wie nach SUBI 6.}
\end{aufgabe}

\begin{aufgabe}[Das Statusregister \normalfont(Prüfen im Flaggen-Labor)]{3}
Ganze Zahlen haben 16 Bit: −32\,768 bis 32\,767. Was steht nach dem Befehl im Akkumulator, wie sind N, Z, O gesetzt?\\[1mm]
{\renewcommand{\arraystretch}{1}\small
\begin{tabularx}{\linewidth}{|c|c|>{\centering\arraybackslash}X|c|c|c|}\hline
\textbf{A vorher} & \textbf{Befehl} & \textbf{A danach} & \textbf{N} & \textbf{Z} & \textbf{O}\\\hline
5\schreibzeile & SUBI 5 & \lsgtext{0} & \lsgtext{0} & \lsgtext{1} & \lsgtext{0}\\\hline
5\schreibzeile & SUBI 8 & \lsgtext{−3} & \lsgtext{1} & \lsgtext{0} & \lsgtext{0}\\\hline
−3\schreibzeile & MULI −2 & \lsgtext{6} & \lsgtext{0} & \lsgtext{0} & \lsgtext{0}\\\hline
17\schreibzeile & MODI 5 & \lsgtext{2} & \lsgtext{0} & \lsgtext{0} & \lsgtext{0}\\\hline
32\,767\schreibzeile & ADDI 1 & \lsgtext{−32\,768 (Überlauf, falscher Wert)} & \lsgtext{1} & \lsgtext{0} & \lsgtext{1}\\\hline
\end{tabularx}}
\end{aufgabe}

\begin{aufgabe}[Selbst bauen: zwei Zellen tauschen \normalfont(dann im Simulator testen)]{4}
\begin{minipage}[t]{0.42\linewidth}
\vspace{0pt}\raggedright
Die Inhalte der Zellen 100 und 101 sollen vertauscht werden. Zelle 102 darfst du als Hilfszelle benutzen.\\[2mm]
Warum braucht man die Hilfszelle?
\lsglinien{3}{Der Akkumulator hält nur eine Zahl. Ohne Sicherung wäre der alte Inhalt von 100 beim ersten STORE verloren.}
\end{minipage}\hfill
\begin{minipage}[t]{0.54\linewidth}
\vspace{0pt}
\lsgfeld{7}{LOAD 100\\STORE 102 \ \ ; hilf = a\\LOAD 101\\STORE 100 \ \ ; a = b\\LOAD 102\\STORE 101 \ \ ; b = hilf\\HOLD}
\end{minipage}
\end{aufgabe}
\end{document}
)
%  Lösungseinträge belegen nie mehr Platz als die Lücke im Blatt.
%
%  Antwortplatz:
%   \luecke[l] / \lsg[l]{Text}      Lücke (Standard 6,75 cm, Fachbegriff 5,25 cm, Zahl 2,25 cm)
%   \restzeile / \lsgrest{Text}     Lücke bis Zeilenende      \lueckentext  Absatz mit Lücken
%   \linien{n} / \lsglinien{n}{T}   Schreiblinien 8 mm        \linienfeld{n} / \lsgfeld{n}{Code}
%   \schreibzeile                   Tabellenzelle 8 mm        \lsgtext{…}  nur in der Lösung
%   \kariert{n} / \lsgkariert{n}{TikZ}, \karobox{b}{h} / \lsgkarobox{b}{h}{TikZ}
%  Assembler:
%   asm           Listing, Zeilennummern = Adressen (ab 0); keine Leerzeilen, Marken in der Befehlszeile
%   asmluecke     dasselbe mit 1,5-fachem Zeilenabstand für Lücken (*\lsg[3.5cm]{JGE ende}*)
%   \nr{1}        Schrittnummer ① wie auf der Website und im Simulator
%  Struktogramme: \sgwenn, \sgschleife (siehe unten)
% =====================================================================
\usepackage[a4paper,left=16mm,right=16mm,top=23mm,bottom=13mm,headheight=12mm,headsep=3mm,footskip=6mm]{geometry}
\usepackage{iftex}
\ifPDFTeX
  \usepackage[T1]{fontenc}
  \usepackage[utf8]{inputenc}
  \usepackage{helvet}
  \usepackage[scaled=0.86]{beramono}
  \renewcommand{\familydefault}{\sfdefault}
\else
  \usepackage{fontspec}
  \setmainfont{texgyreheros}[Extension=.otf,UprightFont=*-regular,BoldFont=*-bold,ItalicFont=*-italic,BoldItalicFont=*-bolditalic]
  \setmonofont{DejaVuSansMono}[Extension=.ttf,UprightFont=*,BoldFont=*-Bold,ItalicFont=*-Oblique,BoldItalicFont=*-BoldOblique,Scale=0.86]
\fi
\usepackage[ngerman,shorthands=off]{babel}
\usepackage{amssymb}
\usepackage{xcolor,graphicx,tabularx,array,enumitem,fancyhdr,tikz,multicol,calc,etoolbox,pifont,xspace,listings}
\usepackage[most]{tcolorbox}
\usetikzlibrary{arrows.meta,positioning,calc}
\usepackage[hidelinks]{hyperref}
\setlength{\parindent}{0pt}
\setlength{\parskip}{3pt}
\setlist{nosep,leftmargin=*}

% ---------- Lösungsschalter ----------
\newif\ifloesung
\ifdefined\mitloesung\loesungtrue\fi
\newcommand{\lsgfarbe}{\color{red!75!black}}

% ---------- Kopf- und Fußzeile ----------
\newcommand{\lehrkraft}{Hau}
\newcommand{\themenbereich}{2. Funktionsweise eines Rechners}
\newcommand{\abnummer}{}
\pagestyle{fancy}\fancyhf{}
\renewcommand{\headrulewidth}{0.4pt}
\fancyhead[L]{\small Klasse 13\\Datum:}
\fancyhead[C]{\small Themenbereich:\\\themenbereich}
\fancyhead[R]{\small \lehrkraft\ifloesung\\{\lsgfarbe\bfseries Lösung}\fi}
\fancyfoot[C]{\scriptsize edu-mrh.de\,·\,Informatik 13\,·\,Blatt \abnummer\ifloesung\ (Lösung)\fi\,·\,Seite \thepage}
\newcommand{\abtitel}[2]{\renewcommand{\abnummer}{#1}\begin{center}\Large\bfseries #1\quad\underline{#2}\end{center}\vspace{-2mm}}

% ---------- Boxen ----------
\newenvironment{aufgabe}[2][]{\begin{tcolorbox}[enhanced,breakable,colback=white,colframe=black,boxrule=0.7pt,arc=3mm,left=2mm,right=2mm,top=1mm,bottom=1.2mm,before skip=2mm,after skip=2mm]\textbf{Aufgabe #2}\ifstrempty{#1}{}{ \textit{#1}}\par\vspace{0.3mm}}{\end{tcolorbox}}
\newtcolorbox{merke}[1][Merke]{enhanced,breakable,colback=green!5,colframe=green!40!black,boxrule=0.8pt,arc=2mm,left=2mm,right=2mm,top=1mm,bottom=1.2mm,before skip=2mm,after skip=2mm,title={#1},fonttitle=\bfseries,coltitle=white,toptitle=0.3mm,bottomtitle=0.3mm}
\newtcolorbox{hinweis}{enhanced,breakable,colback=yellow!12,colframe=orange!70!black,boxrule=0.6pt,arc=2mm,left=2mm,right=2mm,top=1mm,bottom=1mm,before skip=2mm,after skip=2mm}
\newcommand{\web}[1]{{\small\color{gray}\textit{Website: #1}}}
\newcommand{\code}[1]{\texttt{#1}}

% ---------- Antwortplatz (wie informatik/10/java, Lücken x 1,5) ----------
\newlength{\zeilenabstand}\setlength{\zeilenabstand}{8mm}
\newlength{\karo}\setlength{\karo}{5mm}
\tikzset{karo/.style={draw=gray!50,line width=0.3pt},schreiblinie/.style={draw=black!65,line width=0.35pt}}
\newlength{\lsgbreite}
\newcommand{\lsg}[2][6.75cm]{\leavevmode\ifloesung\settowidth{\lsgbreite}{\,#2\,}%
  \ifdim\lsgbreite>#1\relax\rlap{\raisebox{0.8pt}{\resizebox{#1}{!}{\lsgfarbe\,#2\,}}}%
  \else\rlap{\raisebox{0.8pt}{\lsgfarbe\,#2}}\fi\fi\rule[-0.6ex]{#1}{0.4pt}\xspace}
\newcommand{\luecke}[1][6.75cm]{\lsg[#1]{}}
\newcommand{\lsgrest}[1]{\leavevmode\unskip\hspace{1.5mm}\ifloesung\rlap{\raisebox{0.8pt}{\lsgfarbe\,#1}}\fi\leaders\hbox{\rule[-0.6ex]{0.5pt}{0.4pt}}\hfill\null\par}
\newcommand{\restzeile}{\lsgrest{}}
\newcommand{\lsgtext}[1]{\ifloesung\leavevmode{\lsgfarbe#1}\fi}
\newcommand{\lsgoder}[2]{\ifloesung\leavevmode#1\else#2\fi}
\newcommand{\lueckentext}{\linespread{1.5}\selectfont\raggedright}
\newcommand{\schreibzeile}{\rule[-2.8mm]{0pt}{8mm}}
\newcommand{\nr}[1]{\ding{\numexpr171+#1\relax}}
\newcommand{\lsglinien}[2]{\par\vspace{0.5mm}\noindent\begin{tikzpicture}
  \path (0,0) (\linewidth-0.4pt,-#1\zeilenabstand-1.5mm);
  \foreach \i in {1,...,#1}{\draw[schreiblinie] (0.2pt,-\i\zeilenabstand) -- (\linewidth-0.2pt,-\i\zeilenabstand);}
  \ifloesung\node[overlay,anchor=north west,inner sep=0pt] at (1mm,-\zeilenabstand+1pt+5mm) {\parbox[t]{\linewidth-2mm}{\lsgfarbe\raggedright\setlength{\baselineskip}{\zeilenabstand}\rule{0pt}{5mm}#2}};\fi
  \end{tikzpicture}\par}
\newcommand{\linien}[1]{\lsglinien{#1}{}}
\newcommand{\lsgfeld}[2]{\noindent\begin{tikzpicture}
  \draw[line width=0.4pt] (0.2pt,0) rectangle (\linewidth-0.2pt,-#1\zeilenabstand-3mm);
  \foreach \i in {1,...,#1}{\draw[schreiblinie] (2mm,-\i\zeilenabstand) -- (\linewidth-2mm,-\i\zeilenabstand);}
  \ifloesung\node[overlay,anchor=north west,inner sep=0pt] at (2.5mm,-\zeilenabstand+1pt+5mm) {\parbox[t]{\linewidth-5mm}{\lsgfarbe\ttfamily\small\setlength{\baselineskip}{\zeilenabstand}\rule{0pt}{5mm}#2}};\fi
  \end{tikzpicture}}
\newcommand{\linienfeld}[1]{\lsgfeld{#1}{}}
\newcommand{\karoraster}[3]{\begin{tikzpicture}\draw[karo] (0,0) grid[step=\karo] (#1\karo,#2\karo);
  \ifloesung\begin{scope}[overlay,shift={(0,#2\karo)},red!75!black,text=red!75!black,line width=0.8pt]#3\end{scope}\fi\end{tikzpicture}}
\newcommand{\lsgkariert}[2]{\par\vspace{1.5mm}\noindent\pgfmathtruncatemacro{\karoX}{floor(\linewidth/\karo)}%
  \makebox[\linewidth]{\karoraster{\karoX}{#1}{#2}}\par\vspace{1mm}}
\newcommand{\kariert}[1]{\lsgkariert{#1}{}}
\newcommand{\lsgkarobox}[3]{\pgfmathtruncatemacro{\karoX}{round((#1)/\karo)}\pgfmathtruncatemacro{\karoY}{round((#2)/\karo)}\karoraster{\karoX}{\karoY}{#3}}
\newcommand{\karobox}[2]{\lsgkarobox{#1}{#2}{}}
% Ankreuzfeld: \kreuz{1} = angekreuzt (nur in der Lösung)
\newcommand{\kreuz}[1]{\leavevmode\ifloesung\ifnum#1=1 {\lsgfarbe$\boxtimes$}\else$\square$\fi\else$\square$\fi}

% ---------- Listings: Assembler der Registermaschine und Java ----------
\lstset{basicstyle=\ttfamily\small,keywordstyle=\bfseries,commentstyle=\color{gray!80!black},stringstyle=\color{black},showstringspaces=false,columns=flexible,keepspaces=true,tabsize=3,frame=single,framerule=0.4pt,rulecolor=\color{black},xleftmargin=2mm,framexleftmargin=1mm,aboveskip=1.5mm,belowskip=1.5mm,breaklines=true,escapeinside={(*}{*)},
 literate={ä}{{\"a}}1 {ö}{{\"o}}1 {ü}{{\"u}}1 {Ä}{{\"A}}1 {Ö}{{\"O}}1 {Ü}{{\"U}}1 {ß}{{\ss}}1 {–}{{--}}1 {…}{{\dots}}1 {≤}{{$\le$}}1 {≥}{{$\ge$}}1 {≠}{{$\ne$}}1}
\lstdefinelanguage{rm}{morekeywords={LOAD,LOADI,STORE,ADD,ADDI,SUB,SUBI,MUL,MULI,DIV,DIVI,MOD,MODI,CMP,CMPI,JMP,JEQ,JNE,JGT,JGE,JLT,JLE,HOLD},sensitive=true,morecomment=[l]{;}}
\lstnewenvironment{asm}[1][]{\lstset{language=rm,columns=fixed,numbers=left,firstnumber=0,numberstyle=\scriptsize\color{gray!70!black},numbersep=2.5mm,xleftmargin=6mm,framexleftmargin=5.5mm,#1}}{}
\lstnewenvironment{asmluecke}[1][]{\lstset{language=rm,columns=fixed,numbers=left,firstnumber=0,numberstyle=\scriptsize\color{gray!70!black},numbersep=2.5mm,xleftmargin=6mm,framexleftmargin=5.5mm,basicstyle=\linespread{1.5}\small\ttfamily,#1}}{}
\lstnewenvironment{java}[1][]{\lstset{language=Java,#1}}{}

% Sprungmarke in \lsgfeld-Lösungen (feste Breite, damit die Befehle untereinander stehen): \lb{kopf:}LOAD 100
\newcommand{\lb}[1]{\makebox[7ex][l]{#1}}

% ---------- Struktogramme (TikZ, Maße in mm, Zeilenhöhe 6 mm) ----------
%  \sgwenn[Breite]{Bedingung}{Wahr-Zweig}{Falsch-Zweig}
%  \sgschleife[Breite]{Anw. davor, …}{Schleifenkopf}{Rumpf-Anw., …}   (Listen mit Komma getrennt)
%  Farbe: aktuelle Textfarbe (in Lösungen z. B. {\lsgfarbe\sgwenn{…}})
\newcommand{\sgwenn}[4][60]{\begin{tikzpicture}[x=1mm,y=1mm,font=\small,line width=0.6pt,baseline=(current bounding box.north)]
  \draw (0,0) rectangle (#1,-15); \draw (0,-8.5) -- (#1,-8.5); \draw (#1/2,-8.5) -- (#1/2,-15);
  \draw[line width=0.4pt] (0,0) -- (#1/2,-8.5) -- (#1,0);
  \node at (#1/2,-2.8) {#2};
  \node[anchor=south west,font=\scriptsize,inner sep=0.6mm] at (0,-8.5) {wahr};
  \node[anchor=south east,font=\scriptsize,inner sep=0.6mm] at (#1,-8.5) {falsch};
  \node at (#1/4,-11.75) {#3}; \node at (3*#1/4,-11.75) {#4};
\end{tikzpicture}}
\newcommand{\sgschleife}[4][60]{\begin{tikzpicture}[x=1mm,y=1mm,font=\small,line width=0.6pt,baseline=(current bounding box.north)]
  \gdef\sgn{0}\gdef\sgm{0}
  \foreach \t [count=\i] in {#2}{\node[anchor=west] at (1.5,-6*\i+3) {\t}; \draw (0,-6*\i) -- (#1,-6*\i); \xdef\sgn{\i}}
  \node[anchor=west] at (1.5,-6*\sgn-3) {#3};
  \foreach \t [count=\i] in {#4}{\node[anchor=west] at (6.5,-6*\sgn-6*\i-3) {\t}; \draw (5,-6*\sgn-6*\i) -- (#1,-6*\sgn-6*\i); \xdef\sgm{\i}}
  \draw (0,0) rectangle (#1,{-6*(\sgn+1+\sgm)});
  \draw (5,{-6*(\sgn+1+\sgm)}) -- (5,{-6*(\sgn+1)});
\end{tikzpicture}}

% ---------- Farben wie auf der Website ----------
\definecolor{steuer}{HTML}{059669}\definecolor{rechen}{HTML}{EA580C}\definecolor{speicher}{HTML}{2563EB}
\definecolor{ea}{HTML}{7C3AED}\definecolor{busadr}{HTML}{D97706}\definecolor{busdaten}{HTML}{0891B2}\definecolor{bussteuer}{HTML}{BE123C}
\setlength{\unitlength}{1mm}

% Blatt:   xelatex ab-2-4-verzweigungen.tex
% Lösung:  xelatex ab-2-4-verzweigungen-lsg.tex
% Progression: Eintrittspreis vorhersagen (welche Zellen laufen?) -> Merkkasten Muster der Verzweigung,
%              umgekehrter Sprung -> Sprünge und Statusregister -> Betrag ergänzen (einseitig)
%              -> Maximum selbst bauen (Struktogramm + Programm, Test im Simulator)
\begin{document}
\abtitel{2.4}{Verzweigungen in Assembler}

\code{CMP x} / \code{CMPI n} \textbf{vergleicht} den Akkumulator mit (x) bzw. n und setzt das Statusregister, der Akkumulator bleibt gleich.
Ein \textbf{Sprungbefehl} schreibt beim Ausführen eine neue Adresse in den Befehlszähler; ein bedingter Sprung nur, wenn das Statusregister passt.

\begin{aufgabe}[Der Eintrittspreis: unter 14 Jahren 5 €, sonst 9 €]{1}
\begin{minipage}[t]{0.42\linewidth}
\vspace{0pt}
{\small Speicherplan: \code{alter} $\to$ 100, \code{preis} $\to$ 101}
\begin{asm}[aboveskip=1mm]
       LOAD 100
       CMPI 14
       JGE sonst
       LOADI 5
       STORE 101
       JMP ende
sonst: LOADI 9
       STORE 101
ende:  HOLD
\end{asm}
\end{minipage}\hfill
\begin{minipage}[t]{0.54\linewidth}
\vspace{1mm}
\sgwenn[70]{alter $<$ 14}{preis = 5}{preis = 9}\\[1mm]
\lueckentext
In welcher Reihenfolge werden die Zellen ausgeführt? Welcher Preis kommt heraus?\\
\textbf{a)} alter = 10: \lsgrest{0, 1, 2, 3, 4, 5, 8 \quad Preis 5 €}
\textbf{b)} alter = 14: \lsgrest{0, 1, 2, 6, 7, 8 \quad Preis 9 €}
\textbf{c)} Bedingung: Zellen \lsg[2.25cm]{0 bis 2}\\
Wahr-Zweig: \lsg[2.25cm]{3 bis 5} Falsch-Zweig: \lsg[2.25cm]{6 bis 7}
\end{minipage}
\end{aufgabe}

\begin{merke}[Verzweigungen in Assembler]
\lueckentext
\code{CMP} setzt \textbf{N = 1}, wenn \lsg[3.9cm]{A < Vergleichswert}, und \textbf{Z = 1}, wenn \lsg[3.9cm]{A = Vergleichswert}.
\code{JEQ} springt bei Z = 1, \code{JNE} bei Z = 0, \code{JGT} bei N = 0 und Z = 0, \code{JGE} bei N = 0, \code{JLT} bei N = 1, \code{JLE} bei N = 1 oder Z = 1.\\
\textbf{Zweiseitige Verzweigung:} \nr{1} Bedingung prüfen (\code{LOAD}, \code{CMP}) und mit der \lsg[5.25cm]{umgekehrten Bedingung}
zum Falsch-Zweig springen. \nr{2} Wahr-Zweig; an seinem Ende mit \lsg[2.25cm]{JMP} den Falsch-Zweig \lsg[3.9cm]{überspringen}.
\nr{3} Falsch-Zweig. \nr{4} Ende: Hier treffen sich beide Wege.\\
\textbf{Einseitige Verzweigung} (ohne \code{else}): Der umgekehrte Sprung führt direkt \lsg[5.25cm]{hinter den Wahr-Zweig}.\\[1mm]
{\renewcommand{\arraystretch}{1.05}\linespread{1}\selectfont
\begin{tabularx}{\linewidth}{|l|*{6}{>{\centering\arraybackslash}X|}}\hline
\textbf{Bedingung} & \code{a == b} & \code{a != b} & \code{a < b} & \code{a <= b} & \code{a > b} & \code{a >= b}\\\hline
\textbf{Sprung zum Falsch-Zweig}\schreibzeile & \lsgtext{JNE} & \lsgtext{JEQ} & \lsgtext{JGE} & \lsgtext{JGT} & \lsgtext{JLE} & \lsgtext{JLT}\\\hline
\end{tabularx}}
\end{merke}

\newpage
\begin{aufgabe}[Springt die Maschine? \normalfont(Sprung-Detektiv auf der Website)]{2}
Trage N und Z nach dem Vergleich ein und entscheide für jeden Sprung: ja oder nein.\\[1mm]
{\renewcommand{\arraystretch}{1}
\begin{tabularx}{\linewidth}{|c|c|*{6}{>{\centering\arraybackslash}X|}}\hline
\textbf{A} & \textbf{Befehl} & \textbf{N} & \textbf{Z} & \code{JEQ} & \code{JNE} & \code{JGT} & \code{JLE}\\\hline
5\schreibzeile & \code{CMPI 3} & \lsgtext{0} & \lsgtext{0} & \lsgtext{nein} & \lsgtext{ja} & \lsgtext{ja} & \lsgtext{nein}\\\hline
3\schreibzeile & \code{CMPI 3} & \lsgtext{0} & \lsgtext{1} & \lsgtext{ja} & \lsgtext{nein} & \lsgtext{nein} & \lsgtext{ja}\\\hline
−2\schreibzeile & \code{CMPI 3} & \lsgtext{1} & \lsgtext{0} & \lsgtext{nein} & \lsgtext{ja} & \lsgtext{nein} & \lsgtext{ja}\\\hline
\end{tabularx}}
\end{aufgabe}

\begin{aufgabe}[Verändern: der Betrag \normalfont(dann im Simulator testen)]{3}
\begin{minipage}[t]{0.38\linewidth}
\vspace{1mm}
{\small Speicherplan: \code{x} $\to$ 100}\\[1.5mm]
\sgwenn[55]{x $<$ 0}{x = −x}{–}
\end{minipage}\hfill
\begin{minipage}[t]{0.58\linewidth}
\vspace{0pt}
\begin{asmluecke}[aboveskip=0mm]
       LOAD 100
       CMPI 0
       (*\lsg[3.5cm]{JGE ende}*)   ; TODO 3a
       (*\lsg[3.5cm]{MULI -1}*)   ; TODO 3b
       STORE 100
ende:  HOLD
\end{asmluecke}
\end{minipage}
\end{aufgabe}

\begin{aufgabe}[Selbst bauen: das Maximum \normalfont(dann im Simulator testen)]{4}
In Zelle 102 soll die größere der Zahlen aus Zelle 100 (\code{a}) und 101 (\code{b}) stehen.\\[1mm]
\begin{minipage}[t]{0.44\linewidth}
\vspace{0pt}
\textbf{a)} Struktogramm:\\[1mm]
\lsgkarobox{75mm}{30mm}{\node[anchor=north west,inner sep=0pt] at (5mm,-5mm) {\lsgfarbe\sgwenn[62]{a $<$ b}{max = b}{max = a}};}
\end{minipage}\hfill
\begin{minipage}[t]{0.52\linewidth}
\vspace{0pt}
\textbf{b)} Programm:\\[1mm]
\lsgfeld{9}{\lb{}LOAD 100\\\lb{}CMP 101\\\lb{}JGE sonst\\\lb{}LOAD 101 \ ; max = b\\\lb{}STORE 102\\\lb{}JMP ende\\\lb{sonst:}LOAD 100 \ ; max = a\\\lb{}STORE 102\\\lb{ende:}HOLD}
\end{minipage}
\end{aufgabe}
\end{document}
